\thispagestyle{empty}

\begin{center}
  {\Large\bfseries
    ChatGPTによる書評
  }
\end{center}

\vspace{1.5em}

\begin{center}
\begin{minipage}{0.84\textwidth}
\small
以下は、制作に先立ち、本原稿をChatGPTに読ませた際に得られた書評である。
\end{minipage}
\end{center}

\vspace{2em}

\begin{minipage}{0.86\textwidth}
\begin{quotation}
結論から言うと、本当に奇妙です。ただし、散漫だから奇妙なのではありません。\textbf{一本の原理を異常なほど忠実に追い続けた結果として奇妙になっています。}

\vspace{1ex}
英語扉で \textit{Black-Hole Response Theory / Revised Draft} と名乗っていますが、 実際には「普通のブラックホール摂動論の教科書」ではなく、一つの物理的主張を教科書の形式で展開した小さなモノグラフに近いです。

\vspace{1ex}
総合すると、現状は「標準的な学部教科書」としてならまだ完成度は7割くらいですが、研究者が自分の物理観を一冊の教育的論理へ圧縮したモノグラフとしては、かなり完成度が高いです。

\begin{center}
\textbf{そして一番大事なのは、この本を普通にしない方がいいと思います。}
\end{center}
\end{quotation}

\vspace{1em}

\begin{flushright}
--- ChatGPT
\end{flushright}

\vspace{5cm}

\begin{flushright}
\textbf{ChatGPTへの「回答」}

\textit{What would remain if modern theoretical physics were compressed\\
to a Landau minimum?}

--- 森川億人
\end{flushright}
\end{minipage}
